\chapter{OLANDESE VOLANTE}\label{olandese-volante}

\hfill\break

Quando tornai al ponte, Chang si alzò e io mi sedetti sulla poltrona di
comando. Sarebbe stato più utile se avessi capito a cosa serviva uno
qualsiasi di quei pulsanti, in seguito avrei avuto bisogno che Skippy mi
mostrasse le basi. C'era un sacco di gente stretta sul ponte di comando,
era affollato.

«I motori di salto hanno completato la ricarica», disse Chang, indicando
una barra verde lungo tutta la parte inferiore dello schermo principale.

«Grande, grazie. Prima di saltare, la nostra nave, l'altra nave, la
fregata\ldots{} ehm, Skippy, la fregata ha un nome?» Dire ``la nave''
stava cominciando a stufarmi, soprattutto ora che ne avevamo due.

«Il nome kristang della fregata è ``Celestiale Fiore Albale della
Gloriosa Vittoria'', o qualcosa di simile.»

«Mi stai prendendo in giro», esclamai. Avrei detto che le lucertole
dessero alle loro navi nomi come Killer di Criceti o Assassino.

«La casta guerriera kristang è piuttosto devota alla poesia. Vedo che
questo ti sorprende.»

«Poesia?», chiese Chang e ci scambiammo uno sguardo. Nessuno di noi
riusciva a figurarsi i guerrieri irriducibili che avevamo incontrato
sulla stazione spaziale seduti a comporre poesie.

«Ho esaurienti esempi di poesia kristang, ti piacerebbe ascoltarli?»

«No! No, grazie.» È certo come la morte che non ero dell'umore adatto
per la poesia delle lucertole. «Lo chiameremo, ehm, la Fiore per ora.
Non voglio chiamare nulla ``Vittoria'' finché non avremo conseguito
qualcosa di significativo. Tornando alla mia domanda, teniamo la Fiore,
o ce ne sbarazziamo? È danneggiata.»

«Signore», Adams prese la parola, «io dico di tenerla. È danneggiata, ma
sappiamo che funziona e come controllarla, più o meno.» Gli sguardi che
mi rivolse mi dissero che potevo anche essere un alto ufficiale, ma non
ne sapevo molto di combattimento: mai rinunciare a un potenziale
vantaggio. Me l'aveva già ricordato una volta.

«Il sergente Adams ha ragione, colonnello Joe», disse Skippy, «il danno
alla Fiore non ne ha degradato in modo significativo la capacità di
combattimento. Forse ci sono situazioni in cui avere una nave kristang
potrebbe rivelarsi utile.»

«Bene allora, per adesso la teniamo.»

«A proposito di nomi, questa nave come la chiamano i thuranin?», chiese
Thompson.

«Spero non sia un altro nome lungo come la messa cantata, tipo quello
della Fiore», commentò Simms con irritazione.

«I thuranin non danno nomi alle loro navi, ogni nave ha una designazione
numerica», disse Skippy.

«Possiamo darglielo noi, allora», riflettei.

«Oh, andiamo, dobbiamo chiamarla Enterprise», disse Adams con
entusiasmo.

«Non è che tutte le astronavi debbano chiamarsi Enterprise», protestò il
maggiore Simms. Ne dedussi che non fosse una fan di Star Trek.

«Questa è la prima astronave umana. Se non contiamo la nave kristang che
abbiamo rubato. Si deve chiamare Enterprise», insistette Adams.

«L'America non è l'unica cultura con veicoli spaziali immaginari
famosi», protestò Chang, «dovremmo considerare\ldots»

«La nave non è vostra, scimmie!», si inserì Skippy. «Voi siete
l'equipaggio al mio servizio. Se c'è qualcuno che può dare un nome a
questa nave, quello sono io.»

Agitai le braccia per attirare l'attenzione di tutti: «Che ne dite se
prima ci concentriamo su dove stiamo andando e ci preoccupiamo dei nomi
più tardi? Dove andiamo? Come contattiamo il Collettivo, Skippy?».

«Se esiste ancora», iniziò a dire lui.

«Se? Non sei sicuro che esista ancora? Sai dove trovare questo
Collettivo?»

«Se per ``dove'' intendi la galassia della Via Lattea e il suo paio di
galassie nane annesse, sì. Più numerosi ammassi stellari. Oltre a
questo, non tanto.»

«Porca vacca.» Lanciai uno sguardo colpevole alle persone che avevo
trascinato con me alla ricerca del tempo perduto. «Qual è il piano,
allora? Giriamo per la galassia per sempre, cercando di trovare un'altra
intelligenza artificiale degli Anziani? Merda, dovremmo chiamare questa
nave Olandese Volante», borbottai.

«Questo è un nome eccellente, Joe!», disse Skippy con entusiasmo.
«Fatto, ho cambiato l'identificativo di questa nave in Olandese
Volante.»

«Olandese Volante?», chiese Thompson, avviando un inutile giro di
speculazioni da parte dell'equipaggio, tutti a parlarsi addosso l'un
l'altro.

«Si tratta della leggenda di un capitano olandese che uccise un albatro,
gesto che fu la sua maledizione e lo destinò a vagare per mare in
eterno. O qualcosa del genere. La nave non può mai entrare in porto.»

«Pensavo che il tizio che aveva ucciso l'albatro fosse il Vecchio
Marinaio.»

«Ma tu pensa, ovvio che era vecchio se doveva navigare in eterno.»

«Ma l'Olandese è quel totano in quel film di pirati con Johnny Depp?
Cavolo, quel tizio mi faceva venire i brividi.»

«Sì, immagina che alito da pesce marcio aveva quel tizio.»

«No, quello non era Johnny Depp, lui era il pirata con l'eye liner. Il
tizio biondo era destinato a rimanere in mare.»

«Va bene, va bene, basta!», gridai per terminare la conversazione che
altrimenti non sarebbe mai finita. «Skippy», lo guardai con la testa
inclinata, «chiamare questa bagnarola come una nave maledetta non è un
modo per costruire fiducia nell'equipaggio.»

«Stavo solo cercando di rendermi utile», brontolò Skippy.

«E noi non siamo l'equipaggio al tuo servizio. Semmai, siamo pirati.»
Senza l'approvazione dell'Unef, eravamo pirati, fuorilegge.

«Pirati! Joe, questa è la tua allegra banda di pirati! Argh, corpo di
mille balene!» Skippy non era poi tanto male come pirata. «Mi piace!
Un'allegra banda di pirati siamo noi!»

Mi guardai attorno nel compartimento, esaminando le facce. Quella banda
di pirati sarebbe stata tutt'altro che allegra, se Skippy si aspettava
che visitassimo ogni stella della galassia per localizzare il
Collettivo. Avrei dovuto tenere quella conversazione con un gruppo
limitato di persone, non con l'intero equipaggio presente. «Questa
nave\ldots»

«L'Olandese Volante», mi corresse Skippy.

«Bene, l'Olandese», avevo capito che Skippy non avrebbe mollato su
questo punto, «non girerà a vuoto per la galassia fino alla fine dei
tempi, mandando segnali a caso al Collettivo, giusto?»

«Certo che no, scemotto, finiremmo il carburante. Ma tu pensa. Dai
registri dei kristang so dove tengono gli artefatti degli Anziani che
riconosco come relativi a componenti del Collettivo. C'è una struttura
di ricerca in un asteroide, che le lucertole ritengono segreta, a un
paio di wormhole da qui.»

Suonava molto meglio. «Ok, quindi il piano è? Avvicinarci abbastanza
all'asteroide e mandare un segnale a questa roba degli Anziani?»

«Ah! Magari! No, colonnello Joe, non te la caverai così a buon mercato.
No, posso ficcanasare in giro a distanza, ma poi dobbiamo passare al
setaccio il posto per mettere in saccoccia il malloppo di cui ho
bisogno.»

Ficcanasare in giro? Mettere in saccoccia il malloppo? Ma Skippy dove
prendeva il suo slang? «Passare al setaccio? Tipo assalto? Sicuro che i
kristang là siano tutti ricercatori nerd disarmati?»

«Oh, oh, no! S'incazzerebbero come bisce. Questa base sull'asteroide è
dove i kristang cercano di capire come funziona la tecnologia degli
Anziani, quindi la tengono segreta ai thuranin, perché sperano di poter
scavalcare i loro patroni e schiacciarli. I thuranin sanno tutto, ovvio,
non sono così stupidi come pensano le lucertole. Lasciano che conducano
le loro ricerche, hai visto mai che scoprano qualcosa di utile, così i
thuranin potranno piombarci sopra e appropriarsene. Il posto è molto
sorvegliato: reti di rilevamento invisibili, bombe atomiche, laser a
raggi X, tutti gli annessi e connessi. Anche questa nave potrebbe essere
distrutta.»

«È un obiettivo difficile.» Non mi piaceva l'idea di viaggiare su un
bersaglio corazzato. Avevamo colto alla sprovvista l'equipaggio della
Fiore ed eravamo riusciti a prendere quella nave con fatica. Ogni base
che i kristang stessero tentando di tenere segreta ai thuranin sarebbe
stata sempre in allerta. «Dev'esserci un altro posto dove potremmo
andare, o no?»

«No. Non c'è nessun posto adatto. Inoltre, Joe, anche quello che serve a
te è là.»

Cosa mi serviva? Un cheeseburger? «Che cosa?»

«Un modulo di controllo dei wormhole. Il modulo contiene i codici di cui
ho bisogno per chiudere un wormhole.»

Ora ero incazzato: «Skippy, mi hai detto che potevi chiudere quel
wormhole!». Che diavolo c'era ancora che non mi aveva detto?

«Ehi, posso, posso! Per un po'. Quello che posso fare è interrompere la
connessione del wormhole alla rete, ma è solo una cosa temporanea. I
protocolli di rete alla fine ristabilirebbero la connessione e
resetterebbero il wormhole. Per spegnerlo per sempre, mi serve l'intera
serie di codici di controllo, gli Anziani non me l'hanno fornita.»

«Chissà perché», dissi con fare sardonico. «Sei così incredibilmente
affidabile.»

«Lo sono! Non è colpa mia se sei troppo stupido per fare la domanda
giusta. Ricorda, quando fai supposizioni, fai fare la parte del cretino
a te e a me. O a te, in ogni caso.»

Dovetti dar prova di un grande autocontrollo per non calpestare il suo
coperchio lucido sul ponte di comando. Gli Anziani avevano molto di cui
rendere conto, per avere costruito un tale stronzo. A denti stretti,
chiesi con circospezione: «Hai un piano per superare queste griglie di
rilevamento stealth, missili nucleari, laser e astronavi, giusto?».

«Ah, sì», disse con noncuranza, «è un gioco da ragazzi per me. Le loro
stupide griglie possono rilevare tutto quello che vogliono, istruirò i
loro computer principali perché ignorino gli input. Le lucertole non
sapranno che qualcosa è stato rilevato. E le loro armi non sono un
problema, farò spoofing sui loro sistemi di puntamento, in modo che non
possano tenerci sotto tiro. Per un po'. Quando avremo bussato alla
porta, neppure la mia incredibile meraviglia potrà impedire anche alle
più stupide delle lucertole di capire che qualcosa non va.»

«Grande, bene», dissi in tono irritato, «possiamo pianificare come
attraversare quel ponte quando ci arriviamo. Siete tutti pronti per il
salto?»

Chang e Simms si parlavano sopra a vicenda, io cercavo di interrompere
entrambi, e tutti avevano qualcosa da dire. La cosa andò avanti solo per
un secondo o giù di lì, prima che Skippy gridasse non poco, usando
l'interfono della nave: «Basta! Cazzo, siete una truppa di scimmie che
urla in cima agli alberi. Smettetela di blaterare, non riesco neanche a
seguire il filo dei miei pensieri. Fuori! Fuori! Tutti fuori dal ponte
di comando, tranne il colonnello Joe e i piloti!».

Guardai con aria colpevole gli altri, specie Chang e Simms: «Skippy,
possiamo\ldots».

«Colonnello Joe, hai detto che ci servono linee di comunicazione chiare.
Piloteremo questa nave in territorio nemico, attraverso wormhole e,
forse, in combattimento. Avere una truppa di scimmie che mi sbraita
contro non è una comunicazione chiara. Possono stare fuori dal ponte di
comando, nel Cic, ascoltarci e guardare, ma non voglio sentirli, a meno
che tu non apra l'interfono. Sei tu il capitano, sei un colonnello, sei
al comando. Io comunicherò solo con te o con il pilota, mentre questa
nave è in volo.»

Capii cosa intendeva Skippy. Chang e Simms volevano mettere bocca su
tutto, perché non accettavano davvero la mia autorità. Non li biasimavo,
io stesso non riuscivo a crederci. E, a essere sincero, quella fu per me
una grande opportunità per affermare o, a essere più preciso, testare la
mia autorità. Se Chang, Simms, Giraud o chiunque altro non prendevano
sul serio la mia autorità di comando, questo era un ottimo momento per
scoprirlo. «Skippy ha ragione», dichiarai con la voce più profonda e
autorevole che mi usciva, «siamo in troppi qui dentro. L'accesso al
ponte di comando è limitato all'ufficiale di servizio», e in quel
momento l'ufficiale di servizio ero io, «e ai due piloti. Tutti gli
altri, alle postazioni del Cic.» Avevamo bisogno di programmare i turni
di chi sedeva sulla poltrona di comando come ufficiale di servizio,
avevamo anche bisogno di più di due piloti, Desai era stata il nostro
pilota da quando eravamo evasi dalla prigione kristang e doveva essere
stanca ormai, avendo guidato tre veicoli spaziali sconosciuti solo quel
giorno.

Con mio grande sollievo, nessuno oppose obiezioni all'essere bandito dal
ponte di comando, avevano più facile accesso a comandi e schermi nel Cic
comunque, l'area del ponte era troppo affollata.

«La nave è pronta per il salto e la rotta è programmata nel pilota
automatico», annunciò Skippy con un tocco d'impazienza. «Il pilota sa
quale pulsante premere.»

«Pilota?», chiesi.

«Sono pronta, capitano.» Desai aveva un dito posizionato sopra un grande
pulsante sulla console davanti a lei: «Il signor Skippy mi ha guidata
nella programmazione del salto nel pilota automatico, non capisco come
funziona, davvero».

Forse noi scimmie non avremmo mai capito come funzionava il computer di
navigazione del salto.

«Iniziare salto», mi limitai a dire.

E saltammo. Lo schermo sfarfallò e l'unico cambiamento che potei
rilevare consisteva nel fatto che un punto bianco che si trovava
nell'angolo in basso a destra dello schermo era sparito.

«Salto riuscito, capitano», disse Desai, «secondo gli strumenti.»

«Confermo», si limitò a dire Skippy.

«Eccellente», mi rilassai sulla sedia, «qual è il piano ora?»

Ovvio che Skippy aveva una risposta pronta: «Quando i motori di salto
saranno di nuovo carichi, partiremo. Siamo al sicuro, quel salto ci ha
portati abbastanza lontano da rendere del tutto inutili le ricerche
delle navi thuranin».

Guardai lo schermo di stato, che mostrava le navi kristang attaccate
alle piattaforme lungo l'estesa spina dorsale della nave madre. «E i
nostri ospiti indesiderati?»

«Ah, loro. Posso espellere le altre navi kristang e lasciarle qui prima
di saltare. Vuoi ancora tenere la Fiore, giusto?», suggerì Skippy.

Aggrottai le sopracciglia: «Dovremmo tenere la Fiore, sì». Meno male che
avevamo accorciato il nome chilometrico di quella nave. «Le altre navi
non saranno in grado di tornare su Paradiso? O di avvicinarsi abbastanza
da inviare un segnale d'aiuto?»

«C'è un 36\% di probabilità che, concentrando il carburante in una nave
e sacrificando le altre, una singola nave possa fare abbastanza salti
per trovarsi nel raggio di segnalazione utile, sì. La variabile critica
è lo stato di manutenzione dei motori kristang. C'è un 64\% di
probabilità di guasto al motore di salto, per i salti ripetuti.»

«No.»

«No? Le scimmie vogliono vedere i miei calcoli?» La voce di Skippy
suonava divertita.

«No, questa scimmia vuole che ci sia lo 0\% di probabilità che le
lucertole o i thuranin scoprano cosa è successo qui. Se una delle due
specie scopre che la nostra piccola azione piratesca è opera di umani,
la Terra è finita.»

La risposta dell'intelligenza artificiale arrivò con un altro quasi
impercettibile ritardo: «Posso interferire con i computer che
controllano il motore di salto, ma è probabile che sarebbero in grado di
recuperare funzionalità tramite un ripristino da archivi protetti».

«Non è quello che intendevo. Qual è il carico di armi su questa
bagnarola? I thuranin devono avere qualcosa di più forte dei cannoni a
rotaia.»

«Cannoni a rotaia, certo, ma anche missili e maser.»

«Qualcosa che vaporizzi del tutto una nave lucertola? Voglio che non
resti niente, nessuna scatola nera, niente che potrebbe essere
recuperato.» Vedevo che gli altri mi guardavano scettici, chiedendosi
dove stessi andando a parare con le mie domande.

«Una nave madre non è una nave da combattimento; le sue armi sono per lo
più difensive. Su questa nave non ci sono armi che potrebbero
distruggere del tutto quattordici navi kristang, resterebbero rottami
rilevabili. Inoltre, le navi kristang hanno droni che vengono lanciati
in automatico nel caso in cui la nave sia molto danneggiata; questi
droni trasportano i registri di volo e i dati dei sensori della nave e
sono stealth. Ogni nave trasporta più droni, i sensori di questa nave
avrebbero difficoltà a individuare più droni stealth, lanciati da
ciascuna delle quattordici navi kristang, anche se sono io a farli
funzionare. Inoltre, devo avvertire che quelle quattordici navi
kristang, pur con la loro tecnologia minore, rappresentano una minaccia
per questa nave.» Notai che Skippy abbandonava la sua attitudine
saccente quando parlava di uccidere lucertole. «Tutto questo è
accademico, comunque, colonnello Bishop. Come ti ho detto, mi è proibito
usare le armi.»

«Ma a me no.»

Stavolta, la pausa fu abbastanza lunga perché non fossi l'unico a
notarla: «Non vedo come\ldots».

«Tu prepari quelle armi e tieni sotto tiro il bersaglio e io premo il
pulsante, o qualsiasi cosa usino su questa bagnarola», dissi,
guardandomi attorno alla ricerca di qualsiasi cosa potesse somigliare a
un pannello di controllo delle armi. Tutto mi sembrava ancora un incubo
gotico. «Puoi farlo, giusto? La tua programmazione non t'impedisce di
preparare le armi per sparare, basta che non sia tu a iniziare di fatto
la sequenza di fuoco, no?»

Un'altra lunga pausa: «Colonnello, forse ti ho sottovalutato. È quasi
un'idea intelligente». Per la prima volta, avvertii un po' di rispetto
in quella voce artificiale. «Sei sicuro di volerlo fare?»

Attraverso la finestra, potei vedere che il tenente colonnello Chang, il
maggiore Simms e gli altri avevano sentito tutto quello che avevo detto.
«Quando ero in prigione, in attesa di essere giustiziato per essermi
rifiutato di uccidere donne e bambini criceti innocenti, mi sono detto
che queste lucertole naziste potevano andare dritte all'inferno. Ora,
stavo pensando che avremmo potuto limitarci a disabilitare i loro motori
di salto.» Controllai nello schermo la finestra che visualizzava la
poppa e le file di incrociatori e cacciatorpediniere kristang agganciati
alla nave madre. «Ma queste lucertole stanno minacciando la Terra.»
Guardai il maggiore Simms e lei mi fece un cenno con la testa. «Perciò
fottiamoli», dissi, con una rabbia che mi spaventò.

Ci fu una pausa notevole prima che Skippy parlasse: «È possibile che
debba rivedere il mio giudizio su di te».

«Potresti farlo?» Poi mi ricordai dell'uso iper-letterale che Skippy
faceva della lingua inglese. «Lo farai?»

Non ci fu pausa questa volta: «Posso preparare le armi in modo che tu
possa usarle. Sì. Quali armi vuoi che siano attivate?».

Quello era ancora un problema. Nessuna delle armi a bordo della nave
aveva il potere di fare ciò che volevo e, mentre avremmo colpito i
kristang, loro avrebbero potuto rispondere. Se la capacità di salto
della nave madre fosse stata danneggiata, avremmo potuto restare
bloccati nello spazio interstellare per molto tempo. Guardai Chang e
Simms attraverso il vetro: «Qualcuno ha un'idea?».

Prima che una delle due persone nel Cic potesse rispondere, Desai si
girò con la poltrona. «Potremmo saltare da qualche parte davvero lontana
da qualunque stella, così i kristang resterebbero di sicuro indietro?»

«Siamo già nello spazio interstellare profondo. Andare oltre
aumenterebbe la probabilità che i motori di salto dei kristang
falliscano», disse Skippy, «tuttavia, ci sarebbe ancora\ldots»

«Skippy», era il mio turno d'interromperlo, «puoi tenere sotto controllo
i motori di salto di quelle navi kristang?» Desai mi aveva dato un'idea.

«Per un po', sì.»

«E non abbiamo bisogno di mandare squadre a bordo di quelle navi, di
spinottarti?» Se così fosse stato, il mio piano non avrebbe funzionato,
non avremmo potuto assalire più navi da guerra kristang nello stesso
momento.

«No, non questa volta. Le piattaforme thuranin sono in collegamento
fisico con le navi all'attracco, per ridurre le firme spettrali, così mi
sono infiltrato nei computer delle navi kristang poco dopo avere preso
il controllo di questa.»

La domanda successiva era la chiave del mio piano. «E quanto
impiegherebbero i kristang a caricare i motori per un breve salto?»

«Le navi da guerra mantengono sempre una carica minima nei loro motori
di salto, per sfuggire alle imboscate», spiegò Skippy. «Tutte le navi
kristang sono in grado di saltare nell'immediato a corto raggio, anche
ora.»

«Bene. Quanto siamo lontani dal più vicino gigante gassoso?», chiesi.

La bocca di Desai formò una O silenziosa, mentre capiva cosa avevo in
mente. «Sta pensando di saltare vicino a un gigante gassoso, espellere i
kristang e poi far saltare le loro navi dentro il pianeta, prima che
possano reagire?»

«Alt!», fece Skippy con fare derisorio. «Aspettate, scimmie. Non potete
far saltare un'astronave dentro un pianeta, la gravità distorcerebbe il
punto di uscita, quindi\ldots{} ah, ehm, ho capito. Sì, certo. Sarebbe
molto efficace per fare a pezzi quelle navi, soprattutto perché
cercherebbero di emergere nello spazio-tempo già occupato dal pianeta.
Oh, non l'ho mai visto coi miei occhi. Sarà fantastico! Wooow!» Sembrava
allegro. «Sì, c'è un gigante gassoso grande come Saturno a due salti da
qui. Si trova in un sistema stellare disabitato. Pilota, rotta
impostata, pronti quando i motori di salto raggiungono il 64\% di
carica.»

«Aspetta!», mi affrettai a dire, prima che Desai potesse girare la
poltrona. Avrei dovuto fidarmi del fatto che non avrebbe azionato i
comandi senza il mio ordine, perché teneva le mani in aria, non in
bilico sopra i pulsanti. «Skippy, puoi espellere quelle navi, andare via
e farle saltare dentro questo pianeta grande come Saturno, prima che
possano lanciare quei droni?»

«Ti prego, colonnello Joe, piano con le offese. È una passeggiata.»

\hfill\break

Chang non aveva obiezioni o un'idea migliore. Simms era entusiasta e
vidi in lei un nuovo rispetto, e Desai era tutta occhi. Saltammo due
volte attraverso il vuoto dello spazio interstellare e poi, quando Desai
premette un pulsante, in un batter d'occhio l'immagine dello spazio
interstellare nero che compariva negli schermi fu sostituita da un
gigante gassoso grigio-blu. Il pianeta riempiva tutto il visore, voglio
dire, Skippy doveva aver fatto sfoggio delle sue abilità di navigazione.
Di sicuro sembrava che ci fossimo saltati molto vicino. Non ebbi il
tempo di gridare un allarme, perché la nave madre sussultò non appena
quattordici navi kristang e la capsula carica di thuranin addormentati
furono sottoposte con violenza alla separazione di emergenza, e poi la
nostra nave avanzò, salendo sotto la propulsione spazio-normale avviata
da Desai. Skippy non le aveva dato istruzioni di rotta, a parte
andarsene il prima possibile e non indirizzare la nave verso il pianeta.
Potemmo vedere flash quasi immediati da poppa appena le navi kristang
formarono senza volerlo punti di salto. «Pilota, puoi interrompere la
propulsione. Guarda questo», annunciò Skippy con entusiasmo.

«Guardare cosa?»

Il visore fece uno zoom su una sezione delle cime delle nuvole, che
erano all'improvviso illuminate dall'interno. «Ho fatto saltare tutte e
quattordici le navi in un centinaio di chilometri cubi, in modo che i
tratti finali dei loro punti d'ingresso di salto siano sovrapposti, per
essere sicuri che di loro non sia rimasto nulla.»

«Non è un eccesso di potenza distruttiva?»

«L'eccesso di potenza distruttiva è sottovalutato», disse Skippy
autocompiaciuto. «Mmh. Ehm\ldots{} ah!»

«Ehm\ldots{} ah?» La luce all'interno delle nuvole in basso continuava
ad aumentare, assumendo un fulgore virulento, e le loro cime salivano
ribollendo verso di noi. In fretta. «Ehm\ldots{} ah? Che cazzo hai
fatto, Skippy?»

«L'energia rilasciata all'interno del pianeta è stata molto più alta di
quanto mi aspettassi, da qualche parte in alto nell'intervallo dei
petawatt sostenuti. Wow. È diventato autosufficiente. Accidenti.»

«Accidenti?» Le nuvole correvano verso di noi. «Che diavolo è un
petawatt?», gridai.

«Guai. Siamo al sicuro, fuori portata\ldots{} ehm, penso.» Il tono di
voce di Skippy non mi riempì di fiducia.

«Pensi? Desai, portaci lontano dal pianeta, nell'orbita più alta,
comunque la chiami. Saltaci.»

«Sì, sì, signore», disse lei con un grande sorriso in faccia. Era
evidente che pilotare un'astronave gigante era una cosa che le piaceva.
«A tavoletta.»

Vedemmo il pianeta ritrarsi dietro di noi, mentre le cime delle nubi
ribollivano formando in un lampo un fungo davvero enorme, proiettato in
alto sopra l'atmosfera come un bagliore solare. Skippy riferì eccitato
che una parte significativa dell'atmosfera aveva superato la velocità di
fuga e veniva soffiata per sempre nello spazio. Abbastanza perché i
sensori della nave potessero misurare il cambiamento della massa del
pianeta. Il gas atmosferico stava ancora salendo, ma Desai ci fece
sfrecciare più veloci, ero certo che non saremmo stati inghiottiti.
«Skippy, cos'è andato storto?»

«Storto? È stato fantastico! Accidenti! Ho quasi convertito quel pianeta
in una stella minore. Vorrei che avessimo una migliore copertura dei
sensori.»

«``Storto'' significa che è successo qualcosa che non avevi
pianificato», precisai lentamente. Com'era possibile che dovessi
spiegare qualcosa a un'antica macchina super-intelligente?

«Ah, certo, se hai intenzione di fare la punta ai chiodi», disse Skippy
sprezzante. «La prossima volta, farò saltare una nave nemica più in
fondo, verso il nucleo del pianeta. Ma non è divertente, non riusciremmo
a vedere nulla. Sembrerebbe solo un rutto delle nuvole, o qualcosa di
simile, una noia mortale.»

«Signore?», chiese Desai, senza staccare gli occhi dai comandi. «Devo
continuare ad accelerare? Ci stiamo già allontanando dal pianeta a
quindicimila chilometri orari.»

«Eh? Ah, sì, sì, puoi interrompere la propulsione.» Non aveva senso
sprecare energia. Soprattutto perché non avevamo ancora una destinazione
in mente. «Skippy, siamo a posto? Le navi kristang sono tutte distrutte
e hanno portato con sé i droni?»

«Siamo a posto, colonnello Joe. Questo li ha colti del tutto alla
sprovvista, e ho confuso i loro computer in modo che non potessero
reagire in ogni caso. Ora sono una vaga nube di atomi. La stiva di
carico con i thuranin è stata inghiottita e vaporizzata dall'esplosione,
anche loro sono andati.»

Guardai Chang e Simms attraverso il vetro e alzai il pollice. Entrambi
annuirono. «E adesso, Skippy? Ora che ci siamo liberati degli ospiti
indesiderati, stiamo andando ad assaltare la base sull'asteroide, sì?»
Avrei dovuto provare qualcosa dopo aver distrutto kristang e thuranin,
ma non fu così. Nessuna esultanza, nessun senso di colpa, nessuna
soddisfazione, niente. Quelle navi erano solo oggetti, non pensavo agli
esseri senzienti che erano a bordo di esse. Era molto diverso dagli
scontri a fuoco nella boscaglia nigeriana o dalla lotta contro i criceti
su Paradiso. Forse avrei provato qualcosa dopo, dopo un po' di riposo.
In quel momento non m'importava, finché i kristang e i thuranin non
erano più una minaccia.

«Sì. Aspettiamo che i nostri motori di salto si ricarichino. Ho
programmato un lungo salto, quindi abbiamo bisogno di una carica
completa, che richiederà fino a due ore.»

«Desai, ti va di usare le prossime due ore per imparare a pilotare
questa cassa?»

«Sì, sì, signore!»

«Grande. Skippy, programma alcuni punti verso cui lei possa navigare, o
qualsiasi cosa sia previsto fare per l'addestramento di un pilota.
Adesso stiamo per passare attraverso un wormhole, giusto? Quel wormhole
si trova a circa otto anni luce da Paradiso, sulla base di ipotesi che
ho sentito dalla nostra G2. È la nostra intelligence, a livello di
divisione.»

«Ho memorizzato tutti i tuoi acronimi militari, colonnello Joe.»

«Ah, ovvio che l'hai fatto, scusa.» Perché mi prendevo la briga di
spiegare qualcosa a un essere che sapeva tutto quello che c'era in
qualsiasi tipo di dispositivo di archiviazione dati che l'umanità aveva
portato su Paradiso? Più ogni bit di dati che i kristang e i ruhar
avevano sugli esseri umani. Sapeva molto di più sull'umanità di
qualsiasi essere umano. Se lo capisse, nel contesto, era un altro paio
di maniche. «Percorrere otto anni luce richiede circa sedici giorni su
una nave madre thuranin come questa. A meno che, ehm, non ci siano
diversi tipi di nave madre? Non so su quale tipo di mezzo stiamo
viaggiando.»

«Ce ne sono diversi tipi, e i thuranin hanno corazzate, incrociatori e
quelli che si potrebbero chiamare cacciatorpediniere e fregate, anche
navi da trasporto e navi di supporto. Tutte le loro navi madri hanno più
o meno la stessa capacità di salto.»

«Ok, allora, due settimane da qui al wormhole più vicino\ldots»

«Intendi il wormhole che porta a Campo Raggi-X. Che non è il wormhole
più vicino a Paradiso.»

«No?»

«No.» Era esasperante quando Skippy puntualizzava, di solito non
riuscivo a farlo tacere su nessun argomento.

«Verso quale wormhole ci stiamo dirigendo?»

«Quello che si trova a soli cinque anni luce da Paradiso.»

Feci qualche rapido calcolo mentale. «Dieci giorni, allora?»

«Certo, se vuoi andare piano piano, come i thuranin, un saltello alla
volta. Noi procediamo in salti che sono il doppio del normale, e questo
solo perché ho bisogno di calibrare i motori, metterli a punto, come
diresti tu. In un salto potremmo arrivare lontano più del doppio.»

«Non capisco.»

«Questa è la prima cosa intelligente che hai detto da quando ci siamo
incontrati.»

«Bentornato, Skippy, sapevo che lo stronzo era lì dentro da qualche
parte.»

«Farò finta di non aver sentito. La verità, ragazzo mio, è che i
thuranin hanno rubato la tecnologia presumibilmente avanzata del loro
motore di salto molto tempo fa, ma da allora quegli arroganti piccoli
cyborg non hanno fatto alcun progresso nella comprensione del suo
funzionamento. Idioti. Vediamo se trovo un buon modo per spiegarlo: è
come se avessero rubato un'auto e sapessero come metterla in moto e come
ingranare la marcia. Non hanno idea del fatto che il cambio abbia più
marce, ma sanno che l'auto ha un pedale a gas e che il pedale è sul
pavimento. Se ne vanno come lumache in prima, con il motore che urla, e
proseguono strisciando a ritmo dolorosamente lento, ma potrebbero andare
molto più forte. Soprattutto perché i thuranin non sono bravi a copiare
la tecnologia rubata, ed è come se avessero intagliato il cambio da un
blocco di legno massiccio. Idioti. Anche senza modificarla, la loro
copia del cazzo di un motore di salto è capace di prestazioni molto
maggiori, se sono io a controllarla.»

«Non che tu ti stia vantando, o cose simili.»

«No, sono molto modesto. È sorprendente quanto sia modesto, considerando
quanto sono eccezionale.»

«Ah, certo, la tua eccezionale modestia è il tuo tratto più
impressionante.» Su questa strabuzzai gli occhi. «È incredibile quanto
tu sia umile.»

«Sì. Sono molto orgoglioso della mia umiltà. Inoltre, come ha detto uno
di voi umani: ``Non è vantarsi se è vero''.»

«Sì, sì.» Mi stupiva di continuo il fatto che Skippy sapesse così tanto
sulla storia umana. Come può una lucida lattina di birra conoscere
citazioni da Muhammad Ali? «Fare salti più lunghi non danneggerà i
motori? Non voglio restare bloccato qui fuori. Non che la prospettiva di
trascorrere l'eternità su questa nave con te non suoni assolutamente
meravigliosa.»

«Neanch'io voglio trascorrere l'eternità guardando la tua brutta faccia.
No, i thuranin controllano i loro salti forzando troppa energia
attraverso i motori, cosa che li consuma. Io userò un terzo dell'energia
per saltare più del doppio. I motori dureranno molto più a lungo. Ed è
un bene, perché non possiamo portare la nostra nave pirata rubata in una
stazione di riparazione dei thuranin.»

«Noi scimmie possiamo contribuire a mantenere la nave in funzione?
Potresti mostrarci dove si trovano cric e ruota di scorta, nel caso ci
ritroviamo con una gomma a terra.»

«Se restiamo bloccati, potete scendere e spingere.»

Gli feci la linguaccia. È probabile che anche il più intelligente degli
umani non sarebbe stato in grado di aggiustare neppure un gabinetto a
bordo di una nave thuranin; se fosse accaduto qualcosa ai motori o a
qualsiasi altro sistema critico, saremmo risultati del tutto inutili. E,
nel caso in cui Skippy stesse per trascinarci in un lungo viaggio
attraverso la galassia, i bagni erano un sistema critico.

«Ok, allora, stiamo per saltare verso un wormhole. Come lo
attraverseremo? I kristang o i thuranin o anche i ruhar, a questo punto,
non avranno navi a guardia dell'entrata?»

«Come farebbero?»

«Farebbero cosa?»

«La guardia all'entrata.»

«Con, ehm, sai, le navi. O, tipo, una stazione di battaglia o qualcosa
del genere.» Per un secondo, immaginai una Morte Nera²⁶. Solo piena di
lucertole, invece che di truppe d'assalto.

«A cosa servirebbe una stazione di battaglia? Se ne sta lì e basta.»

Non potevo credere che un essere super-intelligente avesse bisogno che
gli spiegassi io che cosa poteva fare una stazione di battaglia
spaziale. Soprattutto perché non ne avevo mai vista una. Lentamente,
aggiunsi: «Se ne sta lì, di fronte all'ingresso, impedendo alle navi non
autorizzate di\ldots».

«Alt!», m'interruppe Skippy, «ho capito il problema. Uhm, davvero non lo
sai? Cazzo, la tua specie è ancora più stupida di quanto pensassi.»

«Cos'è che non so?»

«I wormhole non sono statici. Si spostano di frequente. Una stazione di
battaglia sarebbe del tutto inutile.»

«Oook. No, Signor So-tutto-io, non lo sapevo.» La Bürgermeister non
aveva menzionato quell'importante dettaglio. «Cosa intendi con ``si
spostano di frequente''?»

«Ah, cavolo, ecco che sto per fare una lezione di fisica a un batterio.
Quelli che chiamate wormhole non sono oggetti fisici, sono proiezioni
nello spazio-tempo locale. Le proiezioni saltano intorno, appena sotto
la velocità della luce, in una sorta di schema a forma di otto che copre
circa un anno luce. Una proiezione del wormhole rimane in un posto tra i
diciassette e i novantadue minuti circa, poi si chiude e riappare al
passo successivo lungo il percorso, che potrebbe essere a un quarto di
anno luce di distanza. I wormhole seguono uno schema prestabilito: nel
corso di un ciclo, un wormhole coprirà ogni posizione lungo il percorso,
ma ci sono milioni di posizioni. Il punto è che non esiste che i
kristang, o qualsiasi altra specie al momento in questa galassia,
possano impedire alle navi di usare un wormhole; non è pratico coprire
tutte quelle posizioni. Il mio piano è saltare vicino a dove apparirà un
wormhole e vedere se ci sono altre navi in attesa, so di preciso quando
un wormhole si sposterà nella posizione successiva. Se è libero,
aspettiamo che il wormhole si apra, ci salti davanti, e poi lo
attraversiamo.»

«Uhm.» Avevo immaginato un wormhole come una specie di Stargate, un
grande anello sospeso nello spazio, con un centro incandescente. «Sembra
un buon piano. E non dire: ``Ovvio che lo è''. Perché i wormhole non
stanno in un posto solo?»

«Non ti eri già dato una risposta? Perché se un wormhole restasse in un
posto solo, una specie potrebbe controllarlo, ma tu pensa. E poi perché,
se un wormhole resta aperto più a lungo, ci vuole una quantità
esponenzialmente più grande di energia per sostenere la connessione. E
un wormhole statico alla fine creerebbe una frattura locale nello
spazio-tempo.»

«Un altro buon consiglio di sicurezza, allora.»

«Un eccellente consiglio di sicurezza. È probabile che non sia uno di
quelli di cui voi scimmie dovreste preoccuparvi, nel caso steste
pensando di creare un wormhole con fango e bastoni.»

«Zitto!» Ogni volta che pensavo che Skippy fosse un tipo a posto, mi
diceva qualcosa per ricordarmi che, in fondo, era uno stronzo.

\hfill\break

Dopo esserci allontanati con un salto dal wormhole, concessi una licenza
a tutto l'equipaggio, perché era esausto e Skippy aveva detto che
eravamo al sicuro, avevamo compiuto abbastanza salti da far sì che fosse
molto improbabile che i thuranin ci trovassero. A parte una persona in
infermeria, che poteva riposare sul pavimento, e un ufficiale di
servizio sul lettino del pilota al ponte di comando, volevo che tutti
dormissero otto ore. Eravamo stanchi, eravamo sconvolti dal punto di
vista emotivo e avevamo tutti bisogno di tempo per elaborare quello che
era successo durante una giornata così lunga e movimentata. Per la
maggior parte dell'equipaggio era iniziata quando un Dodo era piovuto
dal cielo nella loro base logistica e quattro umani ne erano usciti e
avevano stordito i ruhar. Per Chang, Simms, Adams e me, era iniziata
prima, quando eravamo scappati dai ruhar e avevamo rubato il Dodo.

Presi il primo turno di servizio, anche se non ero in alcun modo
qualificato come pilota; Skippy aveva programmato due salti di emergenza
e tutto quello che avrei dovuto fare sarebbe stato premere un pulsante
se si fossero presentati dei problemi. Mentre mi sedevo sulla poltrona
del pilota, che era piccola per Desai e decisamente troppo stretta per
me, iniziai a scrivere un resoconto dell'azione sul mio iPad. Un giorno,
speravo, avrei avuto bisogno di fare rapporto alle autorità sulla Terra
di tutto ciò che era successo, ed era meglio buttarlo giù mentre era
ancora fresco nella mia mente. Cazzo, c'era molto da scrivere. Ero solo
a metà, quando Chang mi diede il cambio tre ore dopo e potei tornare
all'alloggiamento a me assegnato, vicino al ponte di comando. Dopo
essermi tolto le scarpe ed essermi accartocciato nel lettino, non
riuscivo a prendere sonno. Piuttosto che starmene lì invano, mi sedetti,
accesi la luce e iniziai a spippolare sull'iPad.

Dagli altoparlanti del mio dispositivo uscì la voce di Skippy:
«Colonnello Joe, dovresti dormire».

«Ho un sacco da fare, Skippy. Per esempio, organizzare un servizio
commemorativo per le quattro persone che abbiamo perso. Non so nemmeno
quale sia la cerimonia corretta per il taoismo o l'induismo.» Mi
accigliai. La verità era che il senso di colpa mi impediva di dormire.

«La tradizione indù prevede la cremazione», Skippy disse la sua in
proposito. «Ho tutte le informazioni possibili e immaginabili sui
rituali religiosi umani.»

«Ah, ehm, grazie, Skippy.» Avrei potuto chiedere a Chang e Desai cosa
fare con i loro compatrioti. Il corpo di Matheson sarebbe rimasto in
magazzino finché non saremmo tornati sulla Terra per una sepoltura
appropriata, a casa. «Penserai che tutto questo sia stupido.»

«Cosa?»

«I rituali religiosi. Sai, la religione stessa.»

«Perché dovrei pensarlo?»

La sua risposta mi sorprese, mi sarei aspettato un commento saccente
sulle scimmie che adorano gli alberi. «Perché gli esseri che ti hanno
creato erano antichi e super-potenti, e, ehm\ldots»

«Joe, Joe, Joe.» Potei visualizzare Skippy che agitava il coperchio
nella mia direzione. «Gli Anziani erano davvero super-potenti; sono
stati in grado di spostare le stelle. Più volte, quando il percorso una
di esse stava per sconvolgere un sistema solare che gli Anziani
pensavano potesse un giorno generare vita intelligente, rimossero la
stella incriminata. Una volta, quando una stella gigante blu minacciava
di diventare una supernova e inondare di radiazioni letali diversi
sistemi solari circostanti, gli Anziani crearono un wormhole e
trasportarono la stella fuori dalla galassia. Fecero saltare di
duecentomila anni luce una stella super-enorme attraverso un wormhole.
Pensaci. La tua domanda presuppone che esseri così potenti, esseri che
avevano davvero raggiunto l'immortalità, che tali esseri, dicevo, non
vedrebbero di buon occhio la religione. Ti sbagli. Di nuovo. Perché
pensi che gli Anziani ci abbiano lasciato, trasceso?»

«Perché», risposi lentamente per darmi il tempo di pensare, «erano
annoiati? Di questa esistenza?»

«No.» Non c'era traccia di umorismo o del consueto sarcasmo nella voce
di Skippy. «È stato perché avevano dato risposta a tutte le domande
fisiche sull'universo e, quando si raggiunge la fine del fisico, ciò che
rimane dev'essere metafisico. Al di là del mondo naturale, si trova solo
il soprannaturale. Gli Anziani sono stati in grado di ritracciare il
tempo fino all'inizio dell'universo e sono rimasti con una domanda: cosa
c'era prima? Da dove viene l'universo? Da dove vengono loro? La fisica,
anche la fisica così avanzata da toccare il regno della magia, ha un
limite. Gli Anziani trascesero perché volevano sapere, dovevano sapere,
da dove venivano. Volevano comunicare con Dio. O con il loro concetto di
Dio.»

«Basta.» Non riuscivo a pensare a nient'altro da dire.

«Fate lo stesso errore che fanno la maggior parte delle giovani specie
quando pensano agli Anziani: li considerate solo esseri mitici, quasi
divini. Erano persone reali, come tutte le altre. Un tempo, erano
stupidi come voi scimmie, ma poi hanno imparato e si sono evoluti in
qualcosa di meraviglioso. Li ricordo come persone sagge, premurose,
potenti, gentili. Mi mancano e cerco il Collettivo per ricollegarmi con
una parte di loro. Ma so che non sono dei.»

«Hai ragione, Skippy, non li avevo pensati come persone.» In realtà, non
avevo pensato molto agli Anziani in generale, poiché di loro si sapeva
così poco. «Che aspetto avevano?»

«Purtroppo, non so dirtelo. È fastidioso. Ho ricordi di loro, in fondo
alla mia mente, ma quando cerco di immaginarli, non ci riesco, non in
modo consapevole. E la mia programmazione non mi permette di dirti nulla
di ciò che so. Quello che m'infastidisce è che non riesco a capire se
dipende dalla mia programmazione originale o è un errore tecnico. Che è
un altro motivo per cui devo contattare il Collettivo.»

«Lo faremo, Skippy, te l'ho promesso, li troveremo.» Avevo le palpebre
così pesanti che non riuscivo a tenerle aperte. «Grazie, Skippy, ora
vado a dormire un po'. Parliamo più tardi.»

\hfill\break

Nei giorni successivi, ci abituammo a una routine. Mentre l'Olandese
Volantesaltava, caricava i motori e saltava di nuovo lì intorno,
sistemammo gli alloggiamenti, allestimmo una sala da pranzo in una stiva
di carico, facemmo sgomberare un'altra stiva per usarla come palestra,
assegnammo un turno di servizio a tutti i membri dell'equipaggio, me
compreso, esplorammo la nave e cercammo d'imparare il più possibile su
di essa. Simms ebbe la brillante idea di portare via le lenzuola dalla
Fiore, in modo da poterle distendere sul pavimento della nostra zona
notte thuranin, per i più alti dell'equipaggio. Quando mi fu offerto un
grande materasso kristang, lo rifiutai, poi ne trovai uno nel mio
alloggiamento vicino al ponte di comando. Non protestai, quel maledetto
letto corto thuranin mi stava distruggendo la schiena.

\hfill\break

La seconda notte in cui andai a dormire nel mio alloggiamento, Skippy mi
disse che c'era una sorpresa in uno degli armadietti. «Merda, Skippy»,
esclamai aprendo l'armadietto. Era pieno, pieno zeppo di tubi di
plastica, ciascuno colmo per circa per un terzo. «Che cazzo è tutta
questa roba?»

«Esperimenti di sapori. C'è un'etichetta stampata sul tubo di ciascuno.»

Guardando più da vicino un tubo, vidi scritto ``Cioccolata \#14'' in
caratteri minuscoli. «Quattordici tipi di cioccolata?»

«Ventidue, in realtà. Come ho detto, il cioccolato è un sapore molto
complesso. Inoltre, ne ho provato di diversi tipi: toffee, caramello,
fragola, curry, jalapeño, cheddar, banana, mela, cannella, salsa e
altri. C'è una lista sul tuo telefono, se t'interessa leggere, cosa che
so che non farai. Non sono sicuro dei sapori della frutta, se devo
essere sincero.»

Stappando con attenzione il tubo, lo annusai. «A quanti di questi tubi
corrisponde un pasto?»

«Dipende dal tuo livello di esercizio. Di regola, per un maschio della
tua taglia, un tubo e mezzo è sufficiente, tre volte al giorno.»

Annusando di nuovo il contenuto di quella sottospecie di cannuccia,
cercai di decidere tra ingoiarlo senza assaggiarlo e assaporarne prima
un sorso. Poiché mi ero offerto assaggiatore volontario per
l'equipaggio, inclinai il tubo, finché un paio di gocce mi scivolarono
in bocca. «Non male, non male. Mi ricorda una bustina di cioccolata
calda Swiss Miss che è rimasta in un capanno di caccia per molti anni,
finché qualcuno non l'ha trovata in fondo a un cassetto. Stantia e
ammuffita, con i piccoli marshmallow trasformati in rocce friabili.»

Skippy rise: «Non avevo capito che fossi un intenditore di melma, Joe».

«Frullato, Skippy, chiamalo frullato.» La mela dava una sgradevole
sensazione in bocca, sia oleosa che gessosa e, in qualche modo,
granulosa allo stesso tempo.

«Prova il cioccolato numero 6, il mio modello di gusto prevede che ti
piacerà di più.»

«Devo prima finire questo. Slurp.»

«Non ce n'è bisogno, abbiamo un sacco di melma a bordo, abbastanza per
anni.»

«Ah, meraviglioso.» Il numero 6 aveva un sapore migliore, lo confrontai
con l'altro alternando i sorsi dai due tubi: «Hai ragione, il 6 sa più
di cioccolato fondente e ha meno retrogusto gessoso. Ed è meno oleoso,
lascia una sensazione più vellutata in bocca».

«Eccellente! Bevi, Joe, mancano solo venti tipi di cioccolato.»

«Ascolta una cosa», non riuscivo ad affrontare il pensiero di altre
venti melme in quel momento. Ciò che volevo fare era trangugiarne un
paio in fretta e farla finita. «Dimmi, le, tipo, sei melme al cioccolato
che pensi abbiano il gusto peggiore. Le persone hanno gusti differenti,
quindi quello che io penso sia il cioccolato migliore potrebbe non
essere il preferito di qualcun altro, ma posso eliminare quelli davvero
cattivi.» L'ultima cosa che volevo era che il primo sorso di melma
potesse spegnere l'entusiasmo di qualcuno per l'idea in generale. «E
identificare quelli cattivi ti aiuterà a calibrare il tuo modello di
gusto, no?»

«Buona idea, Joe», Skippy mi concedeva qualche rara lode, «indovinare
che cos'ha un buon sapore per gli esseri umani è stato in realtà una
sfida interessante per me. E non è facile.»

«Lieto di contribuire a tenere lontana la noia. Skippy, devo dire che
sono davvero colpito, il cioccolato numero 6 è piuttosto buono, potrei
gustarmelo a colazione. Hai fatto un buon lavoro, temevo l'idea della
melma, e ora non è così male. Prima che cambi idea», o prima che
decidessi di non avere più fame, «quale cioccolato pensi sia il
peggiore?»

Le melme, ed è così che l'equipaggio decise di chiamarle, perché eravamo
soldati e siamo fatti così, non divennero popolari, ma non furono
neanche un disastro. Ci arrangiammo e inventammo combinazioni di sapori,
come avevamo mescolato e abbinato pasti pronti per creare portate che
non erano mai state nelle intenzioni dei militari. Una melma al
cioccolato più una melma alla banana dava un cioccolato-banana. Melma al
cocco più melma al curry, mescolata e versata su un pasto pronto di
pollo riscaldato, creava l'approssimazione di un piatto thailandese.
Tutti erano d'accordo sul fatto che le melme al gusto di frutta fossero
le peggiori, anche se la disgustosa fragola, mescolata con l'insipida
banana, aveva in realtà un buon sapore. Questa imprevedibilità delle
preferenze del gusto umano a volte buttava Skippy fuori di testa:
semplicemente non riusciva, pur con tutta la sua impressionante potenza
di elaborazione, a far sì che il suo modello di gusto lavorasse con
precisione affidabile. Fu d'aiuto il fatto che, dopo pochi giorni, fu in
grado di eliminare la sensazione oleosa in bocca e la natura gessosa
della melma, mentre per la consistenza granulosa o scabrosa non si
poteva fare niente, in quanto necessaria alla fibra digestiva, secondo
Skippy. L'equipaggio aveva svariate opportunità di lamentarsi del
``cibo'' a bordo dell'Olandese, cosa che ritenni positiva. La gente
aveva bisogno di qualcosa di cui lamentarsi, lamentarsi era un'attività
che rafforzava i legami e non era un serio colpo al morale. Una melma al
mattino, melme per pranzo e la gente non vedeva l'ora di mangiare cibo
``reale'', come un pasto pronto, per cena. I membri dell'equipaggio,
compresi Chang, Simms, Desai e Giraud, cercarono d'indurmi a mangiare
del cibo vero, ma io mi accontentavo di un cracker o di un biscotto di
tanto in tanto, cercando di dare il buon esempio. Quando i membri
dell'equipaggio videro ``il vecchio'', che, incredibilmente, ero io,
fissato con la melma, non sembrò loro poi così male doversene cibare due
volte al giorno.

Inoltre, stavo aspettando un cheeseburger, e non avrei accettato
sostituti puzzolenti.